\documentclass[11pt]{article}

\usepackage[margin=1in]{geometry}
\usepackage{enumitem}
\usepackage{xcolor}
\usepackage{hyperref}
\usepackage{listings}
\usepackage{amssymb}
\usepackage{array}
\usepackage{longtable}

\hypersetup{
    colorlinks=true,
    linkcolor=blue,
    urlcolor=blue
}

\lstdefinestyle{pythonstyle}{
    language=Python,
    basicstyle=\ttfamily\small,
    backgroundcolor=\color{gray!10},
    frame=single,
    breaklines=true,
    showstringspaces=false,
    keywordstyle=\color{blue},
    stringstyle=\color{teal},
    commentstyle=\color{gray}
}

\newcommand{\checkbox}{$\square$}

\begin{document}

\begin{center}
    {\Large \textbf{Assignment 02: Measuring Runtime with Python}}\\[6pt]
    {\large Google Colab, Python, Sorting, and Google Sheets}\\[12pt]
    Name: \underline{\hspace{2.5in}} \hfill Date: \underline{\hspace{1.5in}}
\end{center}

\vspace{0.15in}

\section*{What Is This Assignment About?}

In this assignment, you will use Python to measure how long a program takes to run.

Computers can do many tasks very quickly, but not all tasks take the same amount of time. Some tasks take longer when the amount of data gets larger.

You will run a program that sorts lists of random numbers. You will test different list sizes, record the runtimes, and then make a graph.

\subsection*{Main Question}

\begin{center}
\textbf{What happens to the runtime when the list size gets larger?}
\end{center}

\subsection*{What You Will Practice}

\begin{itemize}
    \item Opening and using Google Colab
    \item Running Python code
    \item Changing a variable in a program
    \item Measuring runtime using a timer
    \item Recording data from program output
    \item Making a graph in Google Sheets
    \item Submitting a Python \texttt{.py} file in Microsoft Teams
\end{itemize}

\section*{What You Need}

\begin{itemize}
    \item A computer, laptop, Chromebook, or iPad with internet access
    \item Your own personal Gmail email address
    \item The Google Chrome web browser, updated to the current version
    \item Access to Microsoft Teams for Education
    \item Access to Google Sheets
\end{itemize}

\noindent
\textbf{Important:} For this class, please use \textbf{Google Chrome} to open Google Colab. Do not use Safari unless your teacher gives you permission.

\vspace{0.1in}

\noindent
If you are using an iPad, install or update the \textbf{Google Chrome} app before starting this assignment.

\vspace{0.1in}

\noindent
If you do not have a personal Gmail email address, you need to create one before starting this assignment.

\begin{center}
    \url{https://accounts.google.com/signup}
\end{center}

\section*{What You Will Submit}

You will submit your Python file in Microsoft Teams.

\begin{center}
\begin{tabular}{|>{\raggedright}p{2.6in}|>{\raggedright\arraybackslash}p{3.4in}|}
\hline
\textbf{Submit This} & \textbf{Do Not Submit This} \\
\hline
\texttt{RuntimeSorting\_YourName.py} & Any file ending in \texttt{.ipynb} \\
\hline
\end{tabular}
\end{center}

\noindent
Your teacher wants the Python file only. A Python file ends with:

\begin{center}
\texttt{.py}
\end{center}

\noindent
Do \textbf{not} submit the Colab notebook file. A notebook file ends with:

\begin{center}
\texttt{.ipynb}
\end{center}

\section*{Part 1: Open the Starter Notebook}

Follow these steps carefully.

\begin{enumerate}
    \item Open the \textbf{Google Chrome} browser.
    \item Go to the starter notebook:

    \begin{center}
    \url{https://colab.research.google.com/github/bcc-teachermark/M513-2026/blob/main/Assignment_02.ipynb}
    \end{center}

    \item If Google asks you to sign in, sign in using your own personal Gmail email address.
    \item When the notebook opens, click or tap \textbf{File}.
    \item Choose \textbf{Save a copy in Drive}.
    \item Work only on your own saved copy.
\end{enumerate}

\noindent Check your progress:

\begin{itemize}
    \item[\checkbox] I opened Google Chrome.
    \item[\checkbox] I opened the starter Colab notebook.
    \item[\checkbox] I signed in with my own personal Gmail email address.
    \item[\checkbox] I saved my own copy in Google Drive.
\end{itemize}

\section*{Part 2: What Is Google Colab?}

Google Colab is a website where you can write and run Python code in your web browser.

A Colab notebook is made of boxes called \textbf{cells}.

There are two common types of cells:

\begin{center}
\begin{tabular}{|>{\raggedright}p{1.8in}|>{\raggedright\arraybackslash}p{4.5in}|}
\hline
\textbf{Cell Type} & \textbf{What It Does} \\
\hline
Text cell & Used for notes, explanations, titles, and instructions. \\
\hline
Code cell & Used for writing and running Python code. \\
\hline
\end{tabular}
\end{center}

For this assignment, you will mostly run and edit \textbf{code cells}.

\section*{Part 3: Run the Teacher Setup Cell}

The first code cell in the notebook is the \textbf{Teacher Setup} cell.

This cell loads a teacher-supplied timer called \texttt{CodeTimer}. You do not need to understand how \texttt{CodeTimer} works inside. You only need to know how to use it.

\subsection*{What CodeTimer Does}

\texttt{CodeTimer} measures how long part of a program takes to run.

You will use it like this:

\begin{lstlisting}[style=pythonstyle]
timer = CodeTimer("My timing test")

timer.start()

# Code to measure goes here

timer.stop()
\end{lstlisting}

\subsection*{Run the Setup Cell}

\begin{enumerate}
    \item Find the first code cell.
    \item Click or tap the play button on the left side of the cell.
    \item Wait until you see a message that says:

    \begin{center}
    \texttt{CodeTimer is ready to use.}
    \end{center}
\end{enumerate}

\noindent Check your progress:

\begin{itemize}
    \item[\checkbox] I ran the Teacher Setup cell.
    \item[\checkbox] I saw the message that \texttt{CodeTimer} is ready.
\end{itemize}

\section*{Part 4: First Timing Test}

The next code cell is a small first timing test.

It will look similar to this:

\begin{lstlisting}[style=pythonstyle]
from codetimer import CodeTimer

timer = CodeTimer("My first timing test")

timer.start()

student_name = "Phat"
print("Hello T. Mark. I am " + student_name + ".")

timer.stop()
\end{lstlisting}

\subsection*{Your Task}

Change this line:

\begin{lstlisting}[style=pythonstyle]
student_name = "Phat"
\end{lstlisting}

so that it uses your own name.

For example:

\begin{lstlisting}[style=pythonstyle]
student_name = "Boontam"
\end{lstlisting}

Then run the cell.

\subsection*{What You Should See}

You should see a greeting with your name and a runtime measurement.

Example:

\begin{lstlisting}[style=pythonstyle]
Hello T. Mark. I am Boontam.
My first timing test took 123.45 microseconds.
\end{lstlisting}

\noindent Your exact runtime may be different. That is normal.

\noindent Check your progress:

\begin{itemize}
    \item[\checkbox] I changed the student name to my own name.
    \item[\checkbox] I ran the first timing test.
    \item[\checkbox] I saw a runtime measurement.
\end{itemize}

\section*{Part 5: The Sorting Runtime Experiment}

Now you will measure how long it takes Python to sort lists of random numbers.

A \textbf{list} is a group of values. In this assignment, each list contains random integers.

The program will sort the numbers from smallest to largest.

\subsection*{The Experiment Question}

\begin{center}
\textbf{As the list size gets larger, what happens to the runtime?}
\end{center}

\subsection*{The Code You Will Run}

The sorting experiment cell will look similar to this:

\begin{lstlisting}[style=pythonstyle]
from codetimer import CodeTimer
import random


def make_random_list(size):
    numbers = []

    for count in range(size):
        random_number = random.randint(1, 1000000)
        numbers.append(random_number)

    return numbers


def selection_sort(numbers):
    for start_index in range(len(numbers)):
        smallest_index = start_index

        for current_index in range(start_index + 1, len(numbers)):
            if numbers[current_index] < numbers[smallest_index]:
                smallest_index = current_index

        numbers[start_index], numbers[smallest_index] = numbers[smallest_index], numbers[start_index]


sizes = [100, 200, 400, 800, 1600]

results = []

print("Sorting Runtime Experiment")
print("--------------------------")
print()

for list_size in sizes:
    numbers = make_random_list(list_size)

    timer = CodeTimer("Sorting " + str(list_size) + " numbers")

    timer.start()

    selection_sort(numbers)

    elapsed_time = timer.stop()

    results.append([list_size, elapsed_time])

    print()

print("Copy this table into Google Sheets:")
print("List Size\tRuntime in Microseconds")

for row in results:
    print(str(row[0]) + "\t" + str(round(row[1], 2)))
\end{lstlisting}

\subsection*{Important Line}

This line controls the list sizes:

\begin{lstlisting}[style=pythonstyle]
sizes = [100, 200, 400, 800, 1600]
\end{lstlisting}

You will run the experiment once using the original sizes. Then you will change this line and run the experiment again.

\section*{Part 6: First Run}

Run the sorting experiment using the original list sizes:

\begin{center}
\texttt{100, 200, 400, 800, 1600}
\end{center}

Record your results below.

\begin{center}
\begin{tabular}{|c|c|}
\hline
\textbf{List Size} & \textbf{Runtime in Microseconds} \\
\hline
100 & \\
\hline
200 & \\
\hline
400 & \\
\hline
800 & \\
\hline
1600 & \\
\hline
\end{tabular}
\end{center}

\noindent Check your progress:

\begin{itemize}
    \item[\checkbox] I ran the sorting experiment with the original list sizes.
    \item[\checkbox] I recorded the runtime results.
\end{itemize}

\section*{Part 7: Second Run}

Now change the list sizes.

Find this line in the code:

\begin{lstlisting}[style=pythonstyle]
sizes = [100, 200, 400, 800, 1600]
\end{lstlisting}

Change it to:

\begin{lstlisting}[style=pythonstyle]
sizes = [200, 400, 800, 1600, 3200]
\end{lstlisting}

Run the sorting experiment again.

Record your new results below.

\begin{center}
\begin{tabular}{|c|c|}
\hline
\textbf{List Size} & \textbf{Runtime in Microseconds} \\
\hline
200 & \\
\hline
400 & \\
\hline
800 & \\
\hline
1600 & \\
\hline
3200 & \\
\hline
\end{tabular}
\end{center}

\noindent Check your progress:

\begin{itemize}
    \item[\checkbox] I changed the list sizes.
    \item[\checkbox] I ran the sorting experiment a second time.
    \item[\checkbox] I recorded the new runtime results.
\end{itemize}

\section*{Part 8: Third Run - Your Choice}

Now choose your own list sizes.

Your list should have at least five values.

You may use one of these examples:

\begin{lstlisting}[style=pythonstyle]
sizes = [150, 300, 600, 1200, 2400]
\end{lstlisting}

or

\begin{lstlisting}[style=pythonstyle]
sizes = [250, 500, 1000, 2000, 4000]
\end{lstlisting}

or you may create your own.

\subsection*{My List Sizes}

Write the list sizes you chose:

\vspace{0.3in}

\noindent
\underline{\hspace{6.2in}}

\vspace{0.2in}

\subsection*{My Third Run Results}

\begin{center}
\begin{tabular}{|c|c|}
\hline
\textbf{List Size} & \textbf{Runtime in Microseconds} \\
\hline
 & \\
\hline
 & \\
\hline
 & \\
\hline
 & \\
\hline
 & \\
\hline
\end{tabular}
\end{center}

\noindent Check your progress:

\begin{itemize}
    \item[\checkbox] I chose my own list sizes.
    \item[\checkbox] I ran the experiment with my own list sizes.
    \item[\checkbox] I recorded the runtime results.
\end{itemize}

\section*{Part 9: Make a Graph in Google Sheets}

Now you will make a graph of your results.

The graph should show:

\begin{center}
\textbf{List Size vs. Runtime}
\end{center}

\subsection*{Steps}

\begin{enumerate}
    \item Open Google Sheets.
    \item Create a blank spreadsheet.
    \item In cell A1, type:

    \begin{center}
    \texttt{List Size}
    \end{center}

    \item In cell B1, type:

    \begin{center}
    \texttt{Runtime in Microseconds}
    \end{center}

    \item Enter the results from one of your runs.
    \item Highlight your data.
    \item Click or tap \textbf{Insert}.
    \item Click or tap \textbf{Chart}.
    \item Choose a \textbf{line chart} or \textbf{scatter chart}.
    \item Make sure the x-axis is \textbf{List Size}.
    \item Make sure the y-axis is \textbf{Runtime in Microseconds}.
\end{enumerate}

\subsection*{What to Look For}

Look at the shape of the graph.

Ask yourself:

\begin{itemize}
    \item Does the runtime increase as the list size gets larger?
    \item Does the graph look like a straight line?
    \item Does the graph curve upward?
    \item Did doubling the list size exactly double the runtime?
\end{itemize}

\noindent Check your progress:

\begin{itemize}
    \item[\checkbox] I entered my data into Google Sheets.
    \item[\checkbox] I made a graph.
    \item[\checkbox] My graph shows list size and runtime.
\end{itemize}

\section*{Part 10: Reflection Questions}

Answer these questions in complete sentences.

\begin{enumerate}
    \item What happened to the runtime as the list size became larger?

    \vspace{0.5in}

    \item Did doubling the list size exactly double the runtime? Explain your answer.

    \vspace{0.5in}

    \item Which list size took the longest to sort?

    \vspace{0.5in}

    \item Why might your runtime results be different from another student's results?

    \vspace{0.5in}

    \item Look at your graph. Does the runtime grow in a straight line, or does it curve upward?

    \vspace{0.5in}

    \item What is one thing you changed in the Python code yourself?

    \vspace{0.5in}
\end{enumerate}

\section*{Part 11: iPad Tips}

If you are using an iPad, these tips will help.

\begin{itemize}
    \item Use the \textbf{Chrome} app, not Safari.
    \item Turn your iPad sideways into \textbf{landscape mode}. This gives you more space.
    \item Tap inside a code cell before typing.
    \item Tap the play button to run the code.
    \item If the keyboard covers part of the screen, scroll slightly or turn the iPad sideways.
    \item If you have an external keyboard for your iPad, you may use it, but it is not required.
    \item Make sure your iPad is connected to Wi-Fi.
\end{itemize}

\section*{Part 12: Common Problems and Fixes}

\begin{center}
\begin{longtable}{|>{\raggedright}p{2.2in}|>{\raggedright\arraybackslash}p{4.1in}|}
\hline
\textbf{Problem} & \textbf{What to Check} \\
\hline
I am using Safari. & Close Safari and open Google Colab in the Chrome app instead. \\
\hline
I cannot edit the notebook. & Make sure you chose \textbf{File} then \textbf{Save a copy in Drive}. Work on your own copy. \\
\hline
Nothing happens when I type. & Make sure you clicked or tapped inside a code cell. \\
\hline
The play button is missing. & Click or tap inside the code cell. The play button usually appears on the left. \\
\hline
I see an error message. & Check spelling, quotation marks, parentheses, brackets, and commas. Small details matter in programming. \\
\hline
\texttt{CodeTimer} is not found. & Go back and run the Teacher Setup cell again. \\
\hline
My graph looks different from another student's graph. & That is normal. Different computers and Colab sessions may produce different runtimes. Look for the general pattern. \\
\hline
I downloaded the wrong file. & You must download the \texttt{.py} file, not the \texttt{.ipynb} notebook file. \\
\hline
\end{longtable}
\end{center}

\section*{Part 13: Download Your Python File}

When your work is finished, download your Python file from Google Colab.

\begin{enumerate}
    \item In Google Colab, click or tap \textbf{File}.
    \item Click or tap \textbf{Download}.
    \item Choose \textbf{Download .py}.
    \item Wait for the file to download.
\end{enumerate}

\noindent
Your downloaded file should end with:

\begin{center}
\texttt{.py}
\end{center}

\noindent
Rename the file using this format:

\begin{center}
\texttt{RuntimeSorting\_YourName.py}
\end{center}

\noindent
Example:

\begin{center}
\texttt{RuntimeSorting\_Phat.py}
\end{center}

\noindent
\textbf{Important:} Do not submit a file ending in \texttt{.ipynb}.

\noindent Check your progress:

\begin{itemize}
    \item[\checkbox] I clicked or tapped \textbf{File}.
    \item[\checkbox] I clicked or tapped \textbf{Download}.
    \item[\checkbox] I selected \textbf{Download .py}.
    \item[\checkbox] I checked that my file ends with \texttt{.py}.
\end{itemize}

\section*{Part 14: Submit in Microsoft Teams}

Submit your assignment in Microsoft Teams for Education.

\begin{enumerate}
    \item Open \textbf{Microsoft Teams}.
    \item Go to the class team for this course.
    \item Open this assignment.
    \item Choose the option to \textbf{attach} or \textbf{upload} your work.
    \item Attach your downloaded Python file.
    \item Check that the attached file ends with \texttt{.py}.
    \item Submit or turn in the assignment.
\end{enumerate}

\noindent
Do not attach the Colab notebook file.

\begin{center}
\begin{tabular}{|c|c|}
\hline
\textbf{Correct} & \textbf{Incorrect} \\
\hline
\texttt{RuntimeSorting\_YourName.py} & \texttt{Assignment\_02.ipynb} \\
\hline
\end{tabular}
\end{center}

\section*{Final Checklist}

Before you submit, make sure you completed each item.

\begin{itemize}
    \item[\checkbox] I opened the starter notebook in Google Colab.
    \item[\checkbox] I saved a copy in Google Drive.
    \item[\checkbox] I ran the Teacher Setup cell.
    \item[\checkbox] I changed the first timing test to use my own name.
    \item[\checkbox] I ran the sorting experiment with the original list sizes.
    \item[\checkbox] I changed the list sizes and ran the experiment again.
    \item[\checkbox] I chose my own list sizes and ran the experiment a third time.
    \item[\checkbox] I recorded my runtime results.
    \item[\checkbox] I made a graph in Google Sheets.
    \item[\checkbox] I answered the reflection questions.
    \item[\checkbox] I downloaded my work as a \texttt{.py} file.
    \item[\checkbox] I submitted only the \texttt{.py} file in Microsoft Teams.
\end{itemize}

\section*{Teacher Check}

\begin{center}
\begin{tabular}{|>{\raggedright}p{4.3in}|c|}
\hline
\textbf{Requirement} & \textbf{Points} \\
\hline
Correct \texttt{.py} file submitted, not \texttt{.ipynb} & 2 \\
\hline
First timing test runs and uses student's own name & 1 \\
\hline
Sorting experiment runs successfully & 2 \\
\hline
Student tested at least three sets of list sizes & 2 \\
\hline
Runtime results were recorded & 1 \\
\hline
Google Sheets graph was created & 1 \\
\hline
Reflection questions were completed & 1 \\
\hline
\textbf{Total} & \textbf{10} \\
\hline
\end{tabular}
\end{center}

\end{document}